% STAGED for v1.1 (Session 48, 23 Nov 2026). Not \input in main.tex yet.
% Sources: notes/parity-attempts/2026-11-20-lww-orbit-test.md, 2026-11-21-lww-main-term.md (both audited);
% literature check notes/sources-lww26.md (Sessions 47, 48).
\subsection{What the L\"u--Wang--Wu Type~II estimate would also give}\label{subsec:lww-transfer}
The preprint \cite{LWW26} claims $\sum_{p\le x}\mu(p+2)=o(\pi(x))$ (\lbl{Claimed}; we
have not verified it). Write $L=\log x$, fix $K\ge1$ and ranges
$P_i=(\exp(L^{a_i}),\exp(2L^{a_i})]$ with $1/10<a_1<\dots<a_K<1/5$, and let
$\mathcal A$ be the set of products $p_1\cdots p_K$ with $p_i\in P_i$. The
argument of \cite{LWW26} uses the target at $n+2$ in two places only (our
reading of its Sections~2--5):
\begin{itemize}
\item its Type~I sums depend on the target only through its sums over reduced
  residue classes to moduli $x^{2/5+o(1)}$;
\item its Type~II bound (\cite[Prop.~4.2 with Prop.~5.3]{LWW26}) applies
  Cauchy--Schwarz over the cofactor $h$ of $n+2=ah$, $a\in\mathcal A$, and then uses
  only an $L^2$ bound for the coefficient of $h$: the multiplicativity of $\mu$
  enters through $\mu(ah)=(-1)^K\mu(h)$, and the cancellation comes from the
  M\"obius coefficient of the prime variable.
\end{itemize}
The same elimination of the target is noted in \cite[Remark~6.17]{Wan26}.

\begin{proposition}[\lbl{Proved here} in outline, conditional on \cite[Props.~4.2 and 5.3]{LWW26}, which are \lbl{Claimed}]\label{prop:lww-transfer}
Suppose that \cite[Props.~4.2 and 5.3]{LWW26} hold as stated, and let $K$ be
large enough for them. Then there are $\gg_K x/(\log x)^2$ primes $p\in(x,2x]$
with $p+2=al$, $a\in\mathcal A$, $l$ prime. In particular
$P^+(p+2)\ge p\exp(-2K(\log p)^{1/5})$ for infinitely many $p$.
\end{proposition}

\begin{proof}[Sketch]
Apply the decomposition of \cite[Prop.~2.1]{LWW26} to
$\sum_{x<n\le2x}\Lambda(n)G(n+2)$ with
$G(N)=\prod_iV_i^{-1}\sum_{a\in\mathcal A,\,a\mid N}\Lambda(N/a)$,
$V_i=\sum_{p\in P_i}1/p$. By the two points above, the Type~II part is
$O(xL^{-A})$ for every target bounded by a power of $L$, and the Type~I part is
unchanged, up to $O(xL^{-A})$, if $\Lambda(h)$ is replaced by the sieve model
$\kappa_D\mathbf 1_{(h,P(D))=1}$ with $D=L^{B}$ and $\kappa_D=\prod_{p\le D}(1-1/p)^{-1}$:
the two have the same distribution in reduced classes to moduli
$x^{2/5+o(1)}$ on average, by Bombieri--Vinogradov and the fundamental lemma.
For the model the sum is computed directly and equals $2C_2x(1+o(1))$, which is
also the Hardy--Littlewood prediction. Since $G(p+2)\ne0$ forces $p+2=al^j$, prime
powers $j\ge2$ contribute $O(x^{1/2+o(1)})$, and each $p$ carries weight
$\ll_KL^2$, the count follows.
\end{proof}

Unconditionally the record is $P^+(p+2)>p^{0.677}$ \cite{BH98}, and every exponent
$c<1$ is known under the Elliott--Halberstam conjecture (as cited in \cite{Din25});
a cofactor $a=x^{o(1)}$ with exactly $K$ prime factors in prescribed ranges goes
beyond that argument. The proposition is therefore a consistency test of
\cite[Prop.~4.2]{LWW26}, not an error found in it and not a result about twin
primes: twins are the case $K=0$, where the Type~II bound gives nothing. We found no
source that states it (search scope in the project's \texttt{notes/sources-lww26.md};
the citing works \cite{HLWW26} and \cite{Wan26} treat other multiplicative or
rough targets). Full details, with the audits, are in the route cards
\texttt{notes/parity-attempts/2026-11-20-lww-orbit-test.md} and
\texttt{2026-11-21-lww-main-term.md}.

% bib items to add with this subsection:
% \bibitem{Wan26} L.~Wang, \emph{Why expansion methods for Chowla do not reach twin primes: a method-limitation
%   theorem and an explicit parity compiler}, preprint (Zenodo, 10.5281/zenodo.23687596, v1.3), 2026.
% \bibitem{HLWW26} Hu, L\"u, Wang, Wu, \emph{On the higher-order Titchmarsh divisor problem}, arXiv:2611.12756v1, 2026.
%   (first names to be taken from the arXiv page before integration)
% \bibitem{BH98} copy from paper/main.tex.
% \bibitem{Din25} Ding, arXiv:2510.04026v2 (exact entry to be taken from the arXiv page).
